% =============================================================================
%  Notas del Profesor — Sesión 16: Apuntadores, Referencias y Arreglos
%  Programación Avanzada — Universidad Panamericana
%  Compilar: pdflatex sesion16_apuntadores_referencias_arreglos_profesor.tex
% =============================================================================
\documentclass[12pt, oneside]{article}
\usepackage[profesor]{progavanzada}

\begin{document}

\portada{Sesión 16}{Apuntadores, Referencias y Arreglos}{2026}

\tableofcontents
\newpage

% =============================================================================
\section{Objetivos de la sesión}
% =============================================================================

Al finalizar la sesión el alumno será capaz de:

\begin{enumerate}
  \item Explicar qué es una dirección de memoria y cómo obtenerla con \cpp{\&}.
  \item Declarar apuntadores y distinguir entre el valor del apuntador y el
        valor apuntado.
  \item Usar \cpp{*} para leer y escribir a través de un apuntador.
  \item Explicar qué es una referencia y en qué se diferencia de un apuntador.
  \item Aplicar aritmética de apuntadores sobre arreglos.
  \item Reconocer la equivalencia \cpp{arr[i]} $\equiv$ \cpp{*(arr+i)}.
  \item Identificar el acceso fuera de límites como comportamiento indefinido
        y sus implicaciones en seguridad.
\end{enumerate}

\begin{notaprofesor}[Duración estimada]
  100 minutos (sesión doble recomendada). Sugerido:\\
  20 min — modelo de memoria y operador \cpp{\&} \\
  25 min — apuntadores: declaración, \cpp{*}, modificar valores, \cpp{nullptr} \\
  15 min — referencias y comparación con apuntadores \\
  25 min — arreglos, aritmética de apuntadores, equivalencia \\
  15 min — UB, acceso fuera de límites, \cpp{vector::at()}
\end{notaprofesor}

% =============================================================================
\section{Prerrequisitos}
% =============================================================================

\begin{itemize}
  \item Variables, tipos de datos y operadores básicos (sesiones 1--10)
  \item Estructuras de control: \cpp{if}, \cpp{for}, \cpp{while} (sesiones 8--12)
  \item Funciones: declaración, parámetros, valor de retorno (sesiones 13--15)
\end{itemize}

\begin{notaprofesor}
  Si hay alumnos con experiencia en Python, Java o JavaScript, la analogía
  más efectiva es: los apuntadores de C++ son como las \textit{referencias}
  de esos lenguajes, pero sin gestión automática de memoria. La diferencia
  clave es que en C++ la desreferencia es explícita y el programador
  controla cuándo y cómo se accede a la memoria.
\end{notaprofesor}

% =============================================================================
\section{Secuencia de clase}
% =============================================================================

\subsection{Apertura — El modelo de memoria (20 min)}

Dibuja en el pizarrón el modelo de memoria como una tabla de celdas numeradas.
Declara en vivo una variable \cpp{int x = 16;} y pregunta:

\begin{itemize}
  \item ¿Qué necesita el compilador saber para guardar \cpp{x}? (tipo, nombre, valor)
  \item ¿Qué sucede en la RAM cuando se ejecuta esa línea?
\end{itemize}

Luego ejecuta en el cuaderno la primera celda y muestra la dirección hex.
Haz hincapié en que la dirección cambia en cada ejecución (ASLR).

\begin{notaprofesor}[Pregunta de sondeo — Modelo de memoria]
  ``Si tengo \cpp{int a = 5; int b = 5;}, ¿\cpp{\&a == \&b}?''\\
  Respuesta: \textbf{No}. Dos variables distintas siempre tienen direcciones
  distintas, aunque tengan el mismo valor. El único caso en que dos nombres
  comparten dirección es una referencia.
\end{notaprofesor}

\subsection{Apuntadores (25 min)}

Introduce la declaración \cpp{int* p;} y enfatiza que el \texttt{*} en la
declaración es parte del \textbf{tipo}, no del nombre. Muestra la distinción
entre \cpp{int* p, q;} (p es apuntador, q es int) vs \cpp{int *p, *q;}.

\begin{notaprofesor}[Error conceptual muy frecuente]
  Los alumnos confunden el \cpp{*} de declaración con el \cpp{*} de
  desreferencia. Dedica tiempo explícito a esto. Un ejercicio útil: dar
  fragmentos de código y pedir que identifiquen el rol del \cpp{*} en cada
  aparición.
\end{notaprofesor}

Demuestra el ciclo completo:
\begin{enumerate}
  \item Declarar \cpp{int* p;}
  \item Asignar \cpp{p = \&x;}
  \item Leer con \cpp{*p}
  \item Escribir con \cpp{*p = 18;} y mostrar que \cpp{x} cambió
\end{enumerate}

Introduce \cpp{nullptr} y la buena práctica de inicializar siempre los
apuntadores.

\subsection{Referencias (15 min)}

Muestra \cpp{int\& y = x;} y pide a los alumnos que predigan qué imprimirá
\cpp{\&x} y \cpp{\&y} antes de ejecutarlo. La mayoría esperará direcciones
distintas — el resultado igual es el \textit{momento aha} de la referencia.

Usa la tabla comparativa (apuntador vs referencia) para consolidar. Insiste
en que la referencia no puede ser reasignada ni nula.

\subsection{Arreglos (25 min)}

Dibuja el layout de memoria de \cpp{int arr[5] = \{10,20,30,40,50\};} en el
pizarrón, con las direcciones. Muestra que \cpp{arr} imprime la misma
dirección que \cpp{\&arr[0]}.

Explica la aritmética de apuntadores con el diagrama de celdas: sumar 1
avanza \cpp{sizeof(int)} = 4 bytes. Ejecuta las celdas de \cpp{arr+1},
\cpp{arr+2} para que vean los incrementos hex.

Presenta la equivalencia \cpp{arr[i]} $\equiv$ \cpp{*(arr+i)} y demuestra
el \textit{array decay} con \cpp{sizeof(arr)} vs \cpp{sizeof(p)}.

\begin{notaprofesor}[Dato curioso para despertar interés]
  Muestra en vivo que \cpp{arr[2]} y \cpp{2[arr]} producen el mismo resultado.
  Explica que es consecuencia de la conmutatividad de la suma en la aritmética
  de apuntadores. Deja claro que \textbf{jamás se usa en código real} pero
  es útil para entender cómo funciona internamente.
\end{notaprofesor}

\subsection{Comportamiento indefinido y seguridad (15 min)}

Ejecuta la celda de acceso fuera de límites y muestra los resultados
impredecibles. Conecta con el Gusano Morris como ejemplo histórico concreto.

Introduce \cpp{std::vector::at()} como alternativa segura. Pregunta:
``¿Por qué C++ no verifica límites por defecto?'' (rendimiento — dejar que
el alumno razone sobre el trade-off).

% =============================================================================
\section{Errores conceptuales frecuentes}
% =============================================================================

\begin{notaprofesor}[FAQ y errores comunes]

\textbf{1. ``¿Por qué necesito el tipo en un apuntador si ya tengo la dirección?''}\\
La dirección dice \textit{dónde} está el dato, pero no \textit{cuánto} ocupa.
Sin el tipo, el compilador no sabe cuántos bytes leer al desreferenciar.

\vspace{0.6em}
\textbf{2. Confundir \cpp{p} con \cpp{*p}}\\
\cpp{p} es la dirección. \cpp{*p} es el valor en esa dirección. Un ejercicio
útil: dar \cpp{int x=5; int* p=\&x;} y pedir el valor de \cpp{p}, \cpp{*p},
\cpp{\&p}, \cpp{**(\&p)}.

\vspace{0.6em}
\textbf{3. Creer que \cpp{int\& r = x; r = y;} reasigna la referencia}\\
No. \cpp{r = y} copia el \textit{valor} de \cpp{y} en \cpp{x}. La referencia
sigue apuntando a \cpp{x}. Las referencias no se pueden reasignar.

\vspace{0.6em}
\textbf{4. ``El arreglo y el apuntador son lo mismo''}\\
Son similares pero no idénticos. \cpp{sizeof}, no poder asignar a \cpp{arr++},
y el comportamiento con \cpp{\&arr} vs \cpp{\&p} los distinguen.

\vspace{0.6em}
\textbf{5. Pensar que el UB ``siempre crashea''}\\
El comportamiento indefinido puede pasar silenciosamente. Esta es la razón
por la que es tan peligroso: el bug puede estar dormido y manifestarse solo
en ciertas condiciones.

\end{notaprofesor}

% =============================================================================
\section{Conexiones con temas posteriores}
% =============================================================================

\begin{itemize}
  \item \textbf{Sesión 17} — Paso por referencia: aplica directamente
        apuntadores y referencias como mecanismo de paso de parámetros.
  \item \textbf{Sesiones 18--20} — Memoria dinámica: \cpp{new}/\cpp{delete},
        que devuelven y liberan apuntadores.
  \item \textbf{Sesiones 21--25} — Orientación a Objetos: \cpp{this} es un
        apuntador al objeto actual; herencia y polimorfismo usan apuntadores
        a la tabla virtual (\textit{vtable}).
  \item \textbf{C++ moderno} — Smart pointers (\cpp{unique\_ptr},
        \cpp{shared\_ptr}) reemplazan los apuntadores crudos en código de
        producción gestionando la memoria automáticamente.
\end{itemize}

% =============================================================================
\section{Material de la sesión}
% =============================================================================

\begin{itemize}
  \item \textbf{Cuaderno:}
        \texttt{sesiones/sesion16\_apuntadores\_referencias\_arreglos.ipynb}
  \item \textbf{Ejemplos:}
        \texttt{materialInicial/apuntadores/01--07\_apuntadores.cpp}
  \item \textbf{Notas alumno:}
        \texttt{notas\_alumno/pdf/sesion16\_apuntadores\_referencias\_arreglos.pdf}
\end{itemize}

% =============================================================================
\section{Rúbrica de ejercicios}
% =============================================================================

\begin{center}
\begin{tabular}{p{6cm}cc}
  \toprule
  \textbf{Ejercicio / Criterio} & \textbf{Puntos} & \textbf{Evidencia} \\
  \midrule
  \multicolumn{3}{l}{\textit{Ejercicio 1 — Swap con apuntadores}} \\
  \quad Usa paso por apuntador o referencia & 3 & parámetros con \cpp{*} o \cpp{\&} \\
  \quad El swap es correcto                 & 4 & valores intercambiados \\
  \quad No usa \cpp{std::swap}              & 3 & implementación propia \\
  \midrule
  \multicolumn{3}{l}{\textit{Ejercicio 2 — Máximo con aritmética de apuntadores}} \\
  \quad Usa \cpp{*(arr+i)} sin corchetes    & 4 & código sin \cpp{[]} \\
  \quad Resultado correcto                  & 4 & máximo identificado \\
  \quad Recorre exactamente \cpp{n} elementos & 2 & sin UB \\
  \midrule
  \multicolumn{3}{l}{\textit{Ejercicio 3 — Conversión de temperatura con referencia}} \\
  \quad Usa referencia en el parámetro      & 4 & \cpp{double\&} o \cpp{double*} \\
  \quad Fórmula correcta                    & 4 & $(F - 32) \times 5/9$ \\
  \quad Modifica la variable original       & 2 & sin copia \\
  \midrule
  \textbf{Total}                            & \textbf{30} & \\
  \bottomrule
\end{tabular}
\end{center}

\end{document}
